
    \hypertarget{auswertung-verusch-243}{%
\section{Auswertung verusch 243}\label{auswertung-verusch-243}}

\hypertarget{aufgabe-1-import-und-graphische-darstellung-der-messdaten}{%
\subsection{Aufgabe 1: Import und graphische Darstellung der
Messdaten}\label{aufgabe-1-import-und-graphische-darstellung-der-messdaten}}

Hier werden die Messdaten des Multimeters aufgetragen und Betrachtet

    \begin{tcolorbox}[breakable, size=fbox, boxrule=1pt, pad at break*=1mm,colback=cellbackground, colframe=cellborder]
\prompt{In}{incolor}{40}{\boxspacing}
\begin{Verbatim}[commandchars=\\\{\}]
\PY{k+kn}{import}\PY{+w}{ }\PY{n+nn}{matplotlib}\PY{n+nn}{.}\PY{n+nn}{pyplot}\PY{+w}{ }\PY{k}{as}\PY{+w}{ }\PY{n+nn}{plt}

\PY{k+kn}{import}\PY{+w}{ }\PY{n+nn}{numpy}\PY{+w}{ }\PY{k}{as}\PY{+w}{ }\PY{n+nn}{np}

\PY{k+kn}{import}\PY{+w}{ }\PY{n+nn}{numpy}\PY{+w}{ }\PY{k}{as}\PY{+w}{ }\PY{n+nn}{np}
\PY{c+c1}{\PYZsh{} Schöner für import}
\PY{n}{np}\PY{o}{.}\PY{n}{seterr}\PY{p}{(}\PY{n}{over}\PY{o}{=}\PY{l+s+s1}{\PYZsq{}}\PY{l+s+s1}{ignore}\PY{l+s+s1}{\PYZsq{}}\PY{p}{,} \PY{n}{invalid}\PY{o}{=}\PY{l+s+s1}{\PYZsq{}}\PY{l+s+s1}{ignore}\PY{l+s+s1}{\PYZsq{}}\PY{p}{)}




\PY{n}{f}\PY{p}{,} \PY{n}{gain\PYZus{}roh} \PY{o}{=} \PY{n}{np}\PY{o}{.}\PY{n}{loadtxt}\PY{p}{(}\PY{l+s+s1}{\PYZsq{}}\PY{l+s+s1}{Messung Gruppe2.csv}\PY{l+s+s1}{\PYZsq{}}\PY{p}{,} \PY{n}{skiprows}\PY{o}{=}\PY{l+m+mi}{1}\PY{p}{,} \PY{n}{delimiter}\PY{o}{=}\PY{l+s+s1}{\PYZsq{}}\PY{l+s+s1}{,}\PY{l+s+s1}{\PYZsq{}}\PY{p}{,} \PY{n}{unpack}\PY{o}{=}\PY{k+kc}{True}\PY{p}{,} \PY{n}{usecols}\PY{o}{=}\PY{p}{(}\PY{l+m+mi}{0}\PY{p}{,} \PY{l+m+mi}{1}\PY{p}{)}\PY{p}{)}

\PY{n}{D} \PY{o}{=} \PY{l+m+mf}{1e\PYZhy{}3} \PY{c+c1}{\PYZsh{} Verstärkung}


\PY{n}{g} \PY{o}{=} \PY{n}{gain\PYZus{}roh} \PY{o}{/} \PY{n}{D}


\PY{n}{fig}\PY{p}{,} \PY{n}{ax} \PY{o}{=} \PY{n}{plt}\PY{o}{.}\PY{n}{subplots}\PY{p}{(}\PY{n}{figsize}\PY{o}{=}\PY{p}{(}\PY{l+m+mi}{8}\PY{p}{,} \PY{l+m+mi}{6}\PY{p}{)}\PY{p}{)}

\PY{c+c1}{\PYZsh{} Loglog plot}
\PY{n}{ax}\PY{o}{.}\PY{n}{set\PYZus{}xlabel}\PY{p}{(}\PY{l+s+s1}{\PYZsq{}}\PY{l+s+s1}{Frequenz \PYZdl{}f\PYZdl{} [Hz]}\PY{l+s+s1}{\PYZsq{}}\PY{p}{)}
\PY{n}{ax}\PY{o}{.}\PY{n}{set\PYZus{}ylabel}\PY{p}{(}\PY{l+s+s1}{\PYZsq{}}\PY{l+s+s1}{Verstärkung \PYZdl{}g\PYZdl{}}\PY{l+s+s1}{\PYZsq{}}\PY{p}{)}
\PY{n}{ax}\PY{o}{.}\PY{n}{set\PYZus{}title}\PY{p}{(}\PY{l+s+s1}{\PYZsq{}}\PY{l+s+s1}{Frequenzgang}\PY{l+s+s1}{\PYZsq{}}\PY{p}{)}
\PY{n}{ax}\PY{o}{.}\PY{n}{loglog}\PY{p}{(}\PY{n}{f}\PY{p}{,} \PY{n}{g}\PY{p}{,} \PY{n}{linestyle}\PY{o}{=}\PY{l+s+s1}{\PYZsq{}}\PY{l+s+s1}{None}\PY{l+s+s1}{\PYZsq{}}\PY{p}{,} \PY{n}{marker}\PY{o}{=}\PY{l+s+s1}{\PYZsq{}}\PY{l+s+s1}{.}\PY{l+s+s1}{\PYZsq{}}\PY{p}{,} \PY{n}{label}\PY{o}{=}\PY{l+s+s1}{\PYZsq{}}\PY{l+s+s1}{Messwerte}\PY{l+s+s1}{\PYZsq{}}\PY{p}{)}
\PY{n}{ax}\PY{o}{.}\PY{n}{axis}\PY{p}{(}\PY{p}{[}\PY{l+m+mi}{200}\PY{p}{,} \PY{l+m+mf}{2e5}\PY{p}{,} \PY{l+m+mi}{10}\PY{p}{,} \PY{l+m+mf}{1.5E3}\PY{p}{]}\PY{p}{)}
\PY{n}{ax}\PY{o}{.}\PY{n}{grid}\PY{p}{(}\PY{k+kc}{True}\PY{p}{,} \PY{n}{which}\PY{o}{=}\PY{l+s+s1}{\PYZsq{}}\PY{l+s+s1}{both}\PY{l+s+s1}{\PYZsq{}}\PY{p}{,} \PY{n}{ls}\PY{o}{=}\PY{l+s+s1}{\PYZsq{}}\PY{l+s+s1}{\PYZhy{}\PYZhy{}}\PY{l+s+s1}{\PYZsq{}}\PY{p}{)}
\end{Verbatim}
\end{tcolorbox}

    \begin{center}
    \adjustimage{max size={0.9\linewidth}{0.9\paperheight}}{V243_files/V243_1_0.png}
    \end{center}
    { \hspace*{\fill} \\}
    
    \hypertarget{kommentar}{%
\paragraph{Kommentar}\label{kommentar}}

Hier Lässt sich der erwartete Frequenzgang des Bandpassfilters anhand
der symmetrischen Flanken sehr deutlich erkennen. (Der Abgleich stimmt
mit dem gezeigten Plot der Software überein). Es lässt sich ihm
hochfrequenten Bereich kurz vor Beginnn der Flanke bereits ganz leicht
der Frequenzgang des Aufbaus ohne Bandpassfilters erkennen.

    \hypertarget{aufgabe-2-funktion-anpassen}{%
\subsection{Aufgabe 2: Funktion
anpassen}\label{aufgabe-2-funktion-anpassen}}

Nun wird die Bandpassfunktion an die Messdaten gefittet

    \begin{tcolorbox}[breakable, size=fbox, boxrule=1pt, pad at break*=1mm,colback=cellbackground, colframe=cellborder]
\prompt{In}{incolor}{41}{\boxspacing}
\begin{Verbatim}[commandchars=\\\{\}]
\PY{k+kn}{from}\PY{+w}{ }\PY{n+nn}{scipy}\PY{n+nn}{.}\PY{n+nn}{optimize}\PY{+w}{ }\PY{k+kn}{import} \PY{n}{curve\PYZus{}fit}

\PY{c+c1}{\PYZsh{} fitting function (Frequenzgang Bandpass)}

\PY{k}{def}\PY{+w}{ }\PY{n+nf}{fit\PYZus{}func}\PY{p}{(}\PY{o}{*}\PY{n}{args}\PY{p}{)}\PY{p}{:}
    \PY{n}{f}\PY{p}{,} \PY{n}{V}\PY{p}{,} \PY{n}{W1}\PY{p}{,} \PY{n}{W2}\PY{p}{,} \PY{n}{n1}\PY{p}{,} \PY{n}{n2} \PY{o}{=} \PY{n}{args}
    \PY{k}{return} \PY{n}{V}\PY{o}{/}\PY{n}{np}\PY{o}{.}\PY{n}{sqrt}\PY{p}{(}\PY{p}{(}\PY{l+m+mi}{1} \PY{o}{+} \PY{p}{(}\PY{n}{f}\PY{o}{/}\PY{n}{W1}\PY{p}{)}\PY{o}{*}\PY{o}{*}\PY{p}{(}\PY{l+m+mi}{2}\PY{o}{*}\PY{n}{n1}\PY{p}{)} \PY{o}{*} \PY{p}{(}\PY{l+m+mi}{1} \PY{o}{+} \PY{p}{(}\PY{n}{W2}\PY{o}{/}\PY{n}{f}\PY{p}{)}\PY{o}{*}\PY{o}{*}\PY{p}{(}\PY{l+m+mi}{2}\PY{o}{*}\PY{n}{n2}\PY{p}{)}\PY{p}{)}\PY{p}{)}\PY{p}{)}

\PY{n}{p0} \PY{o}{=} \PY{p}{[}\PY{l+m+mf}{1e3}\PY{p}{,} \PY{l+m+mf}{1e3}\PY{p}{,} \PY{l+m+mi}{50000}\PY{p}{,} \PY{l+m+mi}{5}\PY{p}{,} \PY{l+m+mi}{5}\PY{p}{]} 
\PY{c+c1}{\PYZsh{} Verstärkung, untere Grenzfrequenz, obere Grenzfrequenz, Filterordnung unten, Filterordnung oben}

\PY{n}{popt}\PY{p}{,} \PY{n}{pcov} \PY{o}{=} \PY{n}{curve\PYZus{}fit}\PY{p}{(}\PY{n}{fit\PYZus{}func}\PY{p}{,} \PY{n}{f}\PY{p}{,} \PY{n}{g}\PY{p}{,} \PY{n}{p0}\PY{p}{)}

\PY{n+nb}{print}\PY{p}{(}\PY{l+s+sa}{f}\PY{l+s+s1}{\PYZsq{}}\PY{l+s+s1}{Fit Parameter: }\PY{l+s+se}{\PYZbs{}n}\PY{l+s+s1}{V=}\PY{l+s+si}{\PYZob{}}\PY{n}{popt}\PY{p}{[}\PY{l+m+mi}{0}\PY{p}{]}\PY{l+s+si}{:}\PY{l+s+s1}{.3f}\PY{l+s+si}{\PYZcb{}}\PY{l+s+s1}{+/\PYZhy{}}\PY{l+s+si}{\PYZob{}}\PY{n}{np}\PY{o}{.}\PY{n}{sqrt}\PY{p}{(}\PY{n}{pcov}\PY{p}{[}\PY{l+m+mi}{0}\PY{p}{,}\PY{+w}{ }\PY{l+m+mi}{0}\PY{p}{]}\PY{p}{)}\PY{l+s+si}{:}\PY{l+s+s1}{.3f}\PY{l+s+si}{\PYZcb{}}\PY{l+s+s1}{, }\PY{l+s+se}{\PYZbs{}n}\PY{l+s+s1}{W1=}\PY{l+s+si}{\PYZob{}}\PY{n}{popt}\PY{p}{[}\PY{l+m+mi}{1}\PY{p}{]}\PY{l+s+si}{:}\PY{l+s+s1}{.3f}\PY{l+s+si}{\PYZcb{}}\PY{l+s+s1}{+/\PYZhy{}}\PY{l+s+si}{\PYZob{}}\PY{n}{np}\PY{o}{.}\PY{n}{sqrt}\PY{p}{(}\PY{n}{pcov}\PY{p}{[}\PY{l+m+mi}{1}\PY{p}{,}\PY{+w}{ }\PY{l+m+mi}{1}\PY{p}{]}\PY{p}{)}\PY{l+s+si}{:}\PY{l+s+s1}{.3f}\PY{l+s+si}{\PYZcb{}}\PY{l+s+s1}{, }\PY{l+s+se}{\PYZbs{}n}\PY{l+s+s1}{W2=}\PY{l+s+si}{\PYZob{}}\PY{n}{popt}\PY{p}{[}\PY{l+m+mi}{2}\PY{p}{]}\PY{l+s+si}{:}\PY{l+s+s1}{.3f}\PY{l+s+si}{\PYZcb{}}\PY{l+s+s1}{+/\PYZhy{}}\PY{l+s+si}{\PYZob{}}\PY{n}{np}\PY{o}{.}\PY{n}{sqrt}\PY{p}{(}\PY{n}{pcov}\PY{p}{[}\PY{l+m+mi}{2}\PY{p}{,}\PY{+w}{ }\PY{l+m+mi}{2}\PY{p}{]}\PY{p}{)}\PY{l+s+si}{:}\PY{l+s+s1}{.3f}\PY{l+s+si}{\PYZcb{}}\PY{l+s+s1}{, }\PY{l+s+se}{\PYZbs{}n}\PY{l+s+s1}{n1=}\PY{l+s+si}{\PYZob{}}\PY{n}{popt}\PY{p}{[}\PY{l+m+mi}{3}\PY{p}{]}\PY{l+s+si}{:}\PY{l+s+s1}{.3f}\PY{l+s+si}{\PYZcb{}}\PY{l+s+s1}{+/\PYZhy{}}\PY{l+s+si}{\PYZob{}}\PY{n}{np}\PY{o}{.}\PY{n}{sqrt}\PY{p}{(}\PY{n}{pcov}\PY{p}{[}\PY{l+m+mi}{3}\PY{p}{,}\PY{+w}{ }\PY{l+m+mi}{3}\PY{p}{]}\PY{p}{)}\PY{l+s+si}{:}\PY{l+s+s1}{.3f}\PY{l+s+si}{\PYZcb{}}\PY{l+s+s1}{, }\PY{l+s+se}{\PYZbs{}n}\PY{l+s+s1}{n2=}\PY{l+s+si}{\PYZob{}}\PY{n}{popt}\PY{p}{[}\PY{l+m+mi}{4}\PY{p}{]}\PY{l+s+si}{:}\PY{l+s+s1}{.3f}\PY{l+s+si}{\PYZcb{}}\PY{l+s+s1}{+/\PYZhy{}}\PY{l+s+si}{\PYZob{}}\PY{n}{np}\PY{o}{.}\PY{n}{sqrt}\PY{p}{(}\PY{n}{pcov}\PY{p}{[}\PY{l+m+mi}{4}\PY{p}{,}\PY{+w}{ }\PY{l+m+mi}{4}\PY{p}{]}\PY{p}{)}\PY{l+s+si}{:}\PY{l+s+s1}{.3f}\PY{l+s+si}{\PYZcb{}}\PY{l+s+s1}{\PYZsq{}}\PY{p}{)}


\PY{k}{if} \PY{n}{ax}\PY{o}{.}\PY{n}{get\PYZus{}legend}\PY{p}{(}\PY{p}{)}\PY{p}{:}
    \PY{n}{ax}\PY{o}{.}\PY{n}{get\PYZus{}legend}\PY{p}{(}\PY{p}{)}\PY{o}{.}\PY{n}{remove}\PY{p}{(}\PY{p}{)}
    
\PY{c+c1}{\PYZsh{}muss man machen dass sich nicht die fits ansammeln}


\PY{n}{ax}\PY{o}{.}\PY{n}{plot}\PY{p}{(}\PY{n}{f}\PY{p}{,} \PY{n}{fit\PYZus{}func}\PY{p}{(}\PY{n}{f}\PY{p}{,} \PY{o}{*}\PY{n}{popt}\PY{p}{)}\PY{p}{,} \PY{l+s+s1}{\PYZsq{}}\PY{l+s+s1}{r\PYZhy{}}\PY{l+s+s1}{\PYZsq{}}\PY{p}{,} \PY{n}{label}\PY{o}{=}\PY{l+s+s1}{\PYZsq{}}\PY{l+s+s1}{Fit}\PY{l+s+s1}{\PYZsq{}}\PY{p}{)}
\PY{n}{ax}\PY{o}{.}\PY{n}{legend}\PY{p}{(}\PY{p}{)}

\PY{n}{fig}
\PY{c+c1}{\PYZsh{} jens wagner nicht super effizienter Python Programmierer}
\end{Verbatim}
\end{tcolorbox}

    \begin{Verbatim}[commandchars=\\\{\}]
Fit Parameter:
V=1012.122+/-0.221,
W1=46614.034+/-21.372,
W2=6196.180+/-20.283,
n1=4.627+/-0.010,
n2=9.801+/-0.016
    \end{Verbatim}
 
            
\prompt{Out}{outcolor}{41}{}
    
    \begin{center}
    \adjustimage{max size={0.9\linewidth}{0.9\paperheight}}{V243_files/V243_4_1.png}
    \end{center}
    { \hspace*{\fill} \\}
    

    \hypertarget{kommentar}{%
\paragraph{Kommentar}\label{kommentar}}

Die Bandpassfunktion

\[ g(f) = \frac{V}{\sqrt{(1+1/(f/\omega_1)^{2n_1})*(1+(f/\omega_2)^{2n_2})}}\]

wurde mithilfe von Scipy optimize mit diesen Parametern gefittet

V=1012.122+/-0.221, W1=46614.034+/-21.372, W2=6196.180+/-20.283,
n1=4.627+/-0.010, n2=9.801+/-0.016

Die Unsicherheiten haben sich aus der Kovarianzmatrix ergeben.

Es ist zu erknnen dass die Bandpassfunktion den Erwartungen von Raphael
und Jens Wagner entspricht, da die angesetzen p0 Werte nicht stark von
den optimalen Fit-Paramtern abweichen.

Auch hier ist wieder zu erkennen, wie der Frequenzgang ohne
Bandpassfilter einen minimalen Einfluss zeigt. Der optimale Fit weicht
etwas von den Realwerten der Messung ab, da bei hohen Frequenzen die
Verstärkung einen stärkeren Abfall gegenüber der Erwartung zeigt. Man
kann vermuten, dass bei diesen Frequezen das Messgerät bereits die reale
Verstärkung nicht mehr gang genau messen kannn.

    \hypertarget{aufgabe-3-nummerische-integration}{%
\subsection{Aufgabe 3 Nummerische
Integration}\label{aufgabe-3-nummerische-integration}}

Hier wird das Integral des Frequenzgangs bestimmt, um die
Rauschbandbreite \(B\) zu bestimmen.

    \begin{tcolorbox}[breakable, size=fbox, boxrule=1pt, pad at break*=1mm,colback=cellbackground, colframe=cellborder]
\prompt{In}{incolor}{42}{\boxspacing}
\begin{Verbatim}[commandchars=\\\{\}]
\PY{k+kn}{import}\PY{+w}{ }\PY{n+nn}{scipy}\PY{n+nn}{.}\PY{n+nn}{integrate}\PY{+w}{ }\PY{k}{as}\PY{+w}{ }\PY{n+nn}{integrate}

\PY{n}{def\PYZus{}fit\PYZus{}func\PYZus{}squared} \PY{o}{=} \PY{k}{lambda} \PY{o}{*}\PY{n}{args}\PY{p}{:} \PY{n}{fit\PYZus{}func}\PY{p}{(}\PY{o}{*}\PY{n}{args}\PY{p}{)}\PY{o}{*}\PY{o}{*}\PY{l+m+mi}{2} \PY{c+c1}{\PYZsh{} hat das Aura?}

\PY{n}{B}\PY{o}{=}\PY{n}{integrate}\PY{o}{.}\PY{n}{quad}\PY{p}{(}\PY{n}{def\PYZus{}fit\PYZus{}func\PYZus{}squared}\PY{p}{,} \PY{n}{f}\PY{p}{[}\PY{l+m+mi}{0}\PY{p}{]}\PY{p}{,} \PY{n}{f}\PY{p}{[}\PY{o}{\PYZhy{}}\PY{l+m+mi}{1}\PY{p}{]}\PY{p}{,} \PY{n}{args}\PY{o}{=}\PY{n+nb}{tuple}\PY{p}{(}\PY{n}{popt}\PY{p}{)}\PY{p}{)}

\PY{c+c1}{\PYZsh{} f[len(f)] (was auch nicht geht weil Arrays ab 0 zählen: es müsste f[len(f)\PYZhy{}1] sein) ist nicht so cool wie f[\PYZhy{}1]}

\PY{n+nb}{print}\PY{p}{(}\PY{l+s+sa}{f}\PY{l+s+s1}{\PYZsq{}}\PY{l+s+s1}{Das Integral betraegt: }\PY{l+s+si}{\PYZob{}}\PY{n}{B}\PY{p}{[}\PY{l+m+mi}{0}\PY{p}{]}\PY{l+s+si}{:}\PY{l+s+s1}{.4e}\PY{l+s+si}{\PYZcb{}}\PY{l+s+s1}{ +/\PYZhy{} }\PY{l+s+si}{\PYZob{}}\PY{n}{B}\PY{p}{[}\PY{l+m+mi}{1}\PY{p}{]}\PY{l+s+si}{:}\PY{l+s+s1}{.4f}\PY{l+s+si}{\PYZcb{}}\PY{l+s+s1}{\PYZsq{}}\PY{p}{)}
\end{Verbatim}
\end{tcolorbox}

    \begin{Verbatim}[commandchars=\\\{\}]
Das Integral betraegt: 4.7620e+10 +/- 68.1574
    \end{Verbatim}

    \hypertarget{kommentar}{%
\paragraph{Kommentar}\label{kommentar}}

Das Integral des Frequenzgangs g(f) wurde mithilfe von Scipy.integrate
bestimmt. Die Rauschbandbreite \(B\) beträgt
\[B=4.762 x 10^{10} \,\text{Hz}\]

Die Unsicherheit des Integrals mit \(\Delta B \approx 70 \, Hz\) ist um
Faktor 9 kleiner und somit vernachlässigbar klein.

    \hypertarget{aufgabe-4-bestimmung-der-boltzmannkonstante}{%
\subsection{Aufgabe 4 Bestimmung der
Boltzmannkonstante}\label{aufgabe-4-bestimmung-der-boltzmannkonstante}}

Die Boltzmannkonstante wird nun mithilfe der in der Einleitung bennanten
Gleichung
\[k_B = \frac{\langle U_{\text{aus}}^2 \rangle - \langle U_V^2 \rangle}{4 T R B} = \frac{m}{4 T B}\]
besitmmt.

m beschreibt dabei die Steigung \[ m= 4kTB \] welche aus einem linearen
Fit mit \(f(0) = 0\) gewonnen wird

    \hypertarget{fitten-der-gerade-an-an-die-rauschspannungen}{%
\subsubsection{Fitten der Gerade an an die
Rauschspannungen}\label{fitten-der-gerade-an-an-die-rauschspannungen}}

Werte kommen aus dem Messprotokoll (Tabelle 3.2)

Der Statistische Fehler muss zusätzlich noch mit dem Fehler des
Voltmeters von 0.3\% quadraisch addiert werden:

\[\Delta U_{aus} = \sqrt{(U_{aus}*0.003)^2 + \Delta U^2}\] Die Formel
für die Differenz lautet: \[D = U_{aus}^2 - U_V^2  \]

Die zugehörige Fehlerformel ist:

\[\Delta D = \sqrt{(2U_{\text{aus}}\Delta U_{\text{aus}})^2 + (2U_V\Delta U_V)^2}\]

    \begin{tcolorbox}[breakable, size=fbox, boxrule=1pt, pad at break*=1mm,colback=cellbackground, colframe=cellborder]
\prompt{In}{incolor}{51}{\boxspacing}
\begin{Verbatim}[commandchars=\\\{\}]
\PY{n}{R} \PY{o}{=} \PY{n}{np}\PY{o}{.}\PY{n}{array}\PY{p}{(}\PY{p}{[}\PY{l+m+mi}{5}\PY{p}{,} \PY{l+m+mi}{10}\PY{p}{,} \PY{l+m+mi}{15}\PY{p}{,} \PY{l+m+mi}{20}\PY{p}{,} \PY{l+m+mi}{25}\PY{p}{,} \PY{l+m+mi}{30}\PY{p}{]}\PY{p}{)}\PY{o}{*}\PY{l+m+mf}{1e3} 
\PY{n}{U\PYZus{}aus} \PY{o}{=} \PY{n}{np}\PY{o}{.}\PY{n}{array}\PY{p}{(}\PY{p}{[}\PY{l+m+mf}{2.4039}\PY{p}{,} \PY{l+m+mf}{3.1146}\PY{p}{,} \PY{l+m+mf}{3.6898}\PY{p}{,} \PY{l+m+mf}{4.1987}\PY{p}{,} \PY{l+m+mf}{4.6533}\PY{p}{,} \PY{l+m+mf}{5.0669}\PY{p}{]}\PY{p}{)}\PY{o}{*}\PY{l+m+mf}{1e\PYZhy{}3} \PY{c+c1}{\PYZsh{}millivolt}
\PY{n}{fehler\PYZus{}U\PYZus{}stat} \PY{o}{=} \PY{n}{np}\PY{o}{.}\PY{n}{array}\PY{p}{(}\PY{p}{[}\PY{l+m+mf}{0.00088}\PY{p}{,} \PY{l+m+mf}{0.00124}\PY{p}{,} \PY{l+m+mf}{0.00148}\PY{p}{,} \PY{l+m+mf}{0.00167}\PY{p}{,} \PY{l+m+mf}{0.00181}\PY{p}{,} \PY{l+m+mf}{0.00178}\PY{p}{]}\PY{p}{)}\PY{o}{*}\PY{l+m+mf}{1e\PYZhy{}3} \PY{c+c1}{\PYZsh{}millivolt}
\PY{n}{fehler\PYZus{}U\PYZus{}aus} \PY{o}{=} \PY{n}{np}\PY{o}{.}\PY{n}{sqrt}\PY{p}{(}\PY{p}{(}\PY{n}{fehler\PYZus{}U\PYZus{}stat}\PY{o}{*}\PY{l+m+mf}{1e\PYZhy{}3}\PY{p}{)}\PY{o}{*}\PY{o}{*}\PY{l+m+mi}{2} \PY{o}{+} \PY{p}{(}\PY{n}{U\PYZus{}aus}\PY{o}{*}\PY{l+m+mf}{0.003}\PY{p}{)}\PY{o}{*}\PY{o}{*}\PY{l+m+mi}{2}\PY{p}{)} \PY{c+c1}{\PYZsh{}millivolt}

\PY{n}{U\PYZus{}V} \PY{o}{=} \PY{l+m+mf}{1.3718}\PY{o}{*}\PY{l+m+mf}{1e\PYZhy{}3} \PY{c+c1}{\PYZsh{}Eigensrauschen}
\PY{n}{fehler\PYZus{}U\PYZus{}V} \PY{o}{=} \PY{n}{np}\PY{o}{.}\PY{n}{sqrt}\PY{p}{(}\PY{p}{(}\PY{l+m+mf}{0.00048}\PY{o}{*}\PY{l+m+mf}{1e\PYZhy{}2}\PY{p}{)}\PY{o}{*}\PY{o}{*}\PY{l+m+mi}{2} \PY{o}{+} \PY{p}{(}\PY{n}{U\PYZus{}V}\PY{o}{*}\PY{l+m+mf}{0.003}\PY{p}{)}\PY{o}{*}\PY{o}{*}\PY{l+m+mi}{2}\PY{p}{)}\PY{o}{*}\PY{l+m+mf}{1e\PYZhy{}3} \PY{c+c1}{\PYZsh{}Eigensrauschen}

\PY{n}{diff} \PY{o}{=} \PY{n}{U\PYZus{}aus}\PY{o}{*}\PY{o}{*}\PY{l+m+mi}{2} \PY{o}{\PYZhy{}} \PY{n}{U\PYZus{}V}\PY{o}{*}\PY{o}{*}\PY{l+m+mi}{2}
\PY{n}{fehler\PYZus{}diff} \PY{o}{=} \PY{n}{np}\PY{o}{.}\PY{n}{sqrt}\PY{p}{(}\PY{p}{(}\PY{l+m+mi}{2}\PY{o}{*}\PY{n}{U\PYZus{}aus}\PY{o}{*}\PY{n}{fehler\PYZus{}U\PYZus{}aus}\PY{p}{)}\PY{o}{*}\PY{o}{*}\PY{l+m+mi}{2} \PY{o}{+} \PY{p}{(}\PY{l+m+mi}{2}\PY{o}{*}\PY{n}{U\PYZus{}V}\PY{o}{*}\PY{n}{fehler\PYZus{}U\PYZus{}V}\PY{p}{)}\PY{o}{*}\PY{o}{*}\PY{l+m+mi}{2}\PY{p}{)}

\PY{n}{fig2}\PY{p}{,} \PY{n}{ax2} \PY{o}{=} \PY{n}{plt}\PY{o}{.}\PY{n}{subplots}\PY{p}{(}\PY{n}{figsize}\PY{o}{=}\PY{p}{(}\PY{l+m+mi}{8}\PY{p}{,} \PY{l+m+mi}{6}\PY{p}{)}\PY{p}{)}

\PY{n}{ax2}\PY{o}{.}\PY{n}{errorbar}\PY{p}{(}\PY{n}{R}\PY{p}{,} \PY{n}{diff}\PY{p}{,} \PY{n}{yerr}\PY{o}{=}\PY{n}{fehler\PYZus{}diff}\PY{p}{,} \PY{n}{fmt}\PY{o}{=}\PY{l+s+s1}{\PYZsq{}}\PY{l+s+s1}{.}\PY{l+s+s1}{\PYZsq{}}\PY{p}{,} \PY{n}{label}\PY{o}{=}\PY{l+s+s1}{\PYZsq{}}\PY{l+s+s1}{Messwerte}\PY{l+s+s1}{\PYZsq{}}\PY{p}{,} \PY{n}{capsize}\PY{o}{=}\PY{l+m+mi}{3}\PY{p}{)}


\PY{n}{linear} \PY{o}{=} \PY{k}{lambda} \PY{n}{x}\PY{p}{,} \PY{n}{m} \PY{p}{:} \PY{n}{m}\PY{o}{*}\PY{n}{x}\PY{o}{+}\PY{l+m+mi}{0}

\PY{n}{popt2}\PY{p}{,} \PY{n}{pcov2} \PY{o}{=} \PY{n}{curve\PYZus{}fit}\PY{p}{(}\PY{n}{linear}\PY{p}{,} \PY{n}{R}\PY{p}{,} \PY{n}{diff}\PY{p}{,} \PY{n}{sigma}\PY{o}{=}\PY{n}{fehler\PYZus{}diff}\PY{p}{)}
\PY{c+c1}{\PYZsh{} Berechne die skalierten Werte für den e\PYZhy{}12 Exponenten}
\PY{n}{m\PYZus{}scaled} \PY{o}{=} \PY{n}{popt2}\PY{p}{[}\PY{l+m+mi}{0}\PY{p}{]} \PY{o}{/} \PY{l+m+mf}{1e\PYZhy{}12}
\PY{n}{m\PYZus{}err\PYZus{}scaled} \PY{o}{=} \PY{n}{np}\PY{o}{.}\PY{n}{sqrt}\PY{p}{(}\PY{n}{pcov2}\PY{p}{[}\PY{l+m+mi}{0}\PY{p}{,} \PY{l+m+mi}{0}\PY{p}{]}\PY{p}{)} \PY{o}{/} \PY{l+m+mf}{1e\PYZhy{}12}

\PY{c+c1}{\PYZsh{} 1. Print\PYZhy{}Ausgabe anpassen}
\PY{n+nb}{print}\PY{p}{(}\PY{l+s+sa}{f}\PY{l+s+s1}{\PYZsq{}}\PY{l+s+s1}{Fit Parameter: }\PY{l+s+se}{\PYZbs{}n}\PY{l+s+s1}{m=}\PY{l+s+si}{\PYZob{}}\PY{n}{m\PYZus{}scaled}\PY{l+s+si}{:}\PY{l+s+s1}{.3f}\PY{l+s+si}{\PYZcb{}}\PY{l+s+s1}{e\PYZhy{}12 +/\PYZhy{} }\PY{l+s+si}{\PYZob{}}\PY{n}{m\PYZus{}err\PYZus{}scaled}\PY{l+s+si}{:}\PY{l+s+s1}{.3f}\PY{l+s+si}{\PYZcb{}}\PY{l+s+s1}{e\PYZhy{}12}\PY{l+s+s1}{\PYZsq{}}\PY{p}{)}

\PY{c+c1}{\PYZsh{} 2. Plot\PYZhy{}Beschriftung anpassen (mit sauberem KaTeX\PYZhy{}kompatiblem Text für VS Code)}
\PY{n}{x} \PY{o}{=} \PY{n}{np}\PY{o}{.}\PY{n}{linspace}\PY{p}{(}\PY{l+m+mi}{0}\PY{p}{,} \PY{l+m+mf}{35e3}\PY{p}{,} \PY{l+m+mi}{100}\PY{p}{)}
\PY{n}{ax2}\PY{o}{.}\PY{n}{plot}\PY{p}{(}\PY{n}{x}\PY{p}{,} \PY{n}{linear}\PY{p}{(}\PY{n}{x}\PY{p}{,} \PY{o}{*}\PY{n}{popt2}\PY{p}{)}\PY{p}{,} \PY{l+s+s1}{\PYZsq{}}\PY{l+s+s1}{r\PYZhy{}}\PY{l+s+s1}{\PYZsq{}}\PY{p}{,} 
         \PY{n}{label}\PY{o}{=}\PY{l+s+sa}{rf}\PY{l+s+s1}{\PYZsq{}}\PY{l+s+s1}{Fit mit Steigung \PYZdl{}m=(}\PY{l+s+si}{\PYZob{}}\PY{n}{m\PYZus{}scaled}\PY{l+s+si}{:}\PY{l+s+s1}{.3f}\PY{l+s+si}{\PYZcb{}}\PY{l+s+s1}{ }\PY{l+s+s1}{\PYZbs{}}\PY{l+s+s1}{pm }\PY{l+s+si}{\PYZob{}}\PY{n}{m\PYZus{}err\PYZus{}scaled}\PY{l+s+si}{:}\PY{l+s+s1}{.3f}\PY{l+s+si}{\PYZcb{}}\PY{l+s+s1}{) }\PY{l+s+s1}{\PYZbs{}}\PY{l+s+s1}{cdot 10\PYZca{}}\PY{l+s+se}{\PYZob{}\PYZob{}}\PY{l+s+s1}{\PYZhy{}12}\PY{l+s+se}{\PYZcb{}\PYZcb{}}\PY{l+s+s1}{\PYZdl{}}\PY{l+s+s1}{\PYZsq{}}\PY{p}{)}

\PY{n}{ax2}\PY{o}{.}\PY{n}{set\PYZus{}xlabel}\PY{p}{(}\PY{l+s+s1}{\PYZsq{}}\PY{l+s+s1}{Widerstand \PYZdl{}R\PYZdl{} [Ohm]}\PY{l+s+s1}{\PYZsq{}}\PY{p}{)}
\PY{n}{ax2}\PY{o}{.}\PY{n}{set\PYZus{}ylabel}\PY{p}{(}\PY{l+s+s1}{\PYZsq{}}\PY{l+s+s1}{Quadratische Ausgangsspannung \PYZdl{}U\PYZus{}}\PY{l+s+si}{\PYZob{}aus\PYZcb{}}\PY{l+s+s1}{\PYZca{}2 \PYZhy{} U\PYZus{}V\PYZca{}2\PYZdl{} [V\PYZdl{}\PYZca{}2\PYZdl{}]}\PY{l+s+s1}{\PYZsq{}}\PY{p}{)}
\PY{n}{ax2}\PY{o}{.}\PY{n}{set\PYZus{}title}\PY{p}{(}\PY{l+s+s1}{\PYZsq{}}\PY{l+s+s1}{Quadratische Ausgangsspannung gegen Widerstand (ohne Eigensrauschen)}\PY{l+s+s1}{\PYZsq{}}\PY{p}{)}
\PY{n}{ax2}\PY{o}{.}\PY{n}{grid}\PY{p}{(}\PY{k+kc}{True}\PY{p}{,} \PY{n}{which}\PY{o}{=}\PY{l+s+s1}{\PYZsq{}}\PY{l+s+s1}{both}\PY{l+s+s1}{\PYZsq{}}\PY{p}{,} \PY{n}{ls}\PY{o}{=}\PY{l+s+s1}{\PYZsq{}}\PY{l+s+s1}{\PYZhy{}\PYZhy{}}\PY{l+s+s1}{\PYZsq{}}\PY{p}{)}
\PY{n}{ax2}\PY{o}{.}\PY{n}{legend}\PY{p}{(}\PY{p}{)}
\end{Verbatim}
\end{tcolorbox}

    \begin{Verbatim}[commandchars=\\\{\}]
Fit Parameter:
m=786.557e-12 +/- 2.189e-12
    \end{Verbatim}

            \begin{tcolorbox}[breakable, size=fbox, boxrule=.5pt, pad at break*=1mm, opacityfill=0]
\prompt{Out}{outcolor}{51}{\boxspacing}
\begin{Verbatim}[commandchars=\\\{\}]
<matplotlib.legend.Legend at 0x157918050>
\end{Verbatim}
\end{tcolorbox}
        
    \begin{center}
    \adjustimage{max size={0.9\linewidth}{0.9\paperheight}}{V243_files/V243_11_2.png}
    \end{center}
    { \hspace*{\fill} \\}
    
    \hypertarget{kommentar}{%
\paragraph{Kommentar}\label{kommentar}}

Die gefittete Gerade besitzt die Steigung
\(m = (786.6\,\pm 2.2) \cdot 10^{-12} \,\frac{\text{V}^2}{\Omega}\).

Dass sie durch den Ursprung verläuft ist beim Fit festgelegt worden

    \hypertarget{dikussion-der-guxfcte-der-fits}{%
\subsubsection{Dikussion der Güte der
Fits:}\label{dikussion-der-guxfcte-der-fits}}

Die Güte der Fits kann anhand des Chi-Quadrat-Tests getestet werden. Der
Chi-Quadrat-Wert wird berechnet, indem die Summe der quadrierten
Differenzen zwischen den beobachteten Werten und den durch das Modell
vorhergesagten Werten, gewichtet durch die Varianz der Beobachtungen,
genommen wird. Ein niedriger Chi-Quadrat-Wert im Vergleich zu den
Freiheitsgraden deutet auf eine gute Anpassung hin, während ein hoher
Wert auf eine schlechte Anpassung hinweist.

    \begin{tcolorbox}[breakable, size=fbox, boxrule=1pt, pad at break*=1mm,colback=cellbackground, colframe=cellborder]
\prompt{In}{incolor}{53}{\boxspacing}
\begin{Verbatim}[commandchars=\\\{\}]
\PY{k+kn}{from}\PY{+w}{ }\PY{n+nn}{scipy}\PY{n+nn}{.}\PY{n+nn}{stats}\PY{+w}{ }\PY{k+kn}{import} \PY{n}{chi2}

\PY{n}{chisquare} \PY{o}{=} \PY{n}{np}\PY{o}{.}\PY{n}{sum}\PY{p}{(}\PY{p}{(}\PY{p}{(}\PY{n}{linear}\PY{p}{(}\PY{n}{R}\PY{p}{,} \PY{o}{*}\PY{n}{popt2}\PY{p}{)} \PY{o}{\PYZhy{}} \PY{n}{diff}\PY{p}{)}\PY{o}{*}\PY{o}{*}\PY{l+m+mi}{2}\PY{o}{/} \PY{n}{fehler\PYZus{}diff}\PY{o}{*}\PY{o}{*}\PY{l+m+mi}{2}\PY{p}{)}\PY{p}{)}
\PY{n}{dof} \PY{o}{=} \PY{l+m+mi}{5} 
\PY{n}{chi2\PYZus{}reduced} \PY{o}{=} \PY{n}{chisquare} \PY{o}{/} \PY{n}{dof}
\PY{n+nb}{print}\PY{p}{(}\PY{l+s+sa}{f}\PY{l+s+s1}{\PYZsq{}}\PY{l+s+s1}{Chi\PYZhy{}Quadrat: }\PY{l+s+si}{\PYZob{}}\PY{n}{chisquare}\PY{l+s+si}{:}\PY{l+s+s1}{.3f}\PY{l+s+si}{\PYZcb{}}\PY{l+s+s1}{, Freiheitsgrade: }\PY{l+s+si}{\PYZob{}}\PY{n}{dof}\PY{l+s+si}{\PYZcb{}}\PY{l+s+s1}{, Reduziertes Chi\PYZhy{}Quadrat: }\PY{l+s+si}{\PYZob{}}\PY{n}{chi2\PYZus{}reduced}\PY{l+s+si}{:}\PY{l+s+s1}{.3f}\PY{l+s+si}{\PYZcb{}}\PY{l+s+s1}{\PYZsq{}}\PY{p}{)}

\PY{n}{prob} \PY{o}{=} \PY{n+nb}{round}\PY{p}{(}\PY{l+m+mi}{1} \PY{o}{\PYZhy{}} \PY{n}{chi2}\PY{o}{.}\PY{n}{cdf}\PY{p}{(}\PY{n}{chisquare}\PY{p}{,} \PY{n}{dof}\PY{p}{)}\PY{p}{,} \PY{l+m+mi}{2}\PY{p}{)}\PY{o}{*}\PY{l+m+mi}{100}

\PY{n+nb}{print}\PY{p}{(}\PY{l+s+sa}{f}\PY{l+s+s1}{\PYZsq{}}\PY{l+s+s1}{Fitwahrscheinlichkeit = }\PY{l+s+si}{\PYZob{}}\PY{n}{prob}\PY{l+s+si}{:}\PY{l+s+s1}{.2f}\PY{l+s+si}{\PYZcb{}}\PY{l+s+s1}{\PYZpc{}}\PY{l+s+s1}{\PYZsq{}}\PY{p}{)}
\end{Verbatim}
\end{tcolorbox}

    \begin{Verbatim}[commandchars=\\\{\}]
Chi-Quadrat: 4.667, Freiheitsgrade: 5, Reduziertes Chi-Quadrat: 0.933
Fitwahrscheinlichkeit = 46.00\%
    \end{Verbatim}

    \hypertarget{kommentar}{%
\paragraph{Kommentar}\label{kommentar}}

Der berechnete reduzierte Chi-Quadrat-Wert beträgt
\(\chi^2_{\text{red}} = 0.933\) und liegt damit sehr nah am Optimalwert
von 1. Es lässt sich darauf schlussfolgern (mit 46\%
Wahrscheinlichkeit), dass die Abweichungen tatsächlich von statistischer
Natur sind, und nur ein sehr kleiner systematischer Fehler vorliegt.

    \hypertarget{bestimmung-der-boltzmannkonstante}{%
\subsubsection{Bestimmung der
Boltzmannkonstante}\label{bestimmung-der-boltzmannkonstante}}

Nun wird mithilfe der Steigung des Fits die Boltzmannkonstante bestimmt.

Nach Anleitung wird ein systematischer Fehle für B von 2\% angenommen.

Die Raumtemperatur beträgt \(T = (22.30\,\pm 0.10)^\circ\text{C}\)

    \begin{tcolorbox}[breakable, size=fbox, boxrule=1pt, pad at break*=1mm,colback=cellbackground, colframe=cellborder]
\prompt{In}{incolor}{54}{\boxspacing}
\begin{Verbatim}[commandchars=\\\{\}]
\PY{k+kn}{import}\PY{+w}{ }\PY{n+nn}{uncertainties}\PY{+w}{ }\PY{k}{as}\PY{+w}{ }\PY{n+nn}{unc}

\PY{n}{m} \PY{o}{=} \PY{n}{unc}\PY{o}{.}\PY{n}{ufloat}\PY{p}{(}\PY{n}{popt2}\PY{p}{[}\PY{l+m+mi}{0}\PY{p}{]}\PY{p}{,} \PY{n}{np}\PY{o}{.}\PY{n}{sqrt}\PY{p}{(}\PY{n}{pcov2}\PY{p}{[}\PY{l+m+mi}{0}\PY{p}{,} \PY{l+m+mi}{0}\PY{p}{]}\PY{p}{)}\PY{p}{)}
\PY{n}{T} \PY{o}{=} \PY{n}{unc}\PY{o}{.}\PY{n}{ufloat}\PY{p}{(}\PY{l+m+mf}{22.3}\PY{o}{+}\PY{l+m+mf}{273.15}\PY{p}{,} \PY{l+m+mf}{0.1}\PY{p}{)} \PY{c+c1}{\PYZsh{} Raumtemperatur in Kelvin}
\PY{n}{B\PYZus{}ufloat} \PY{o}{=} \PY{n}{unc}\PY{o}{.}\PY{n}{ufloat}\PY{p}{(}\PY{n}{B}\PY{p}{[}\PY{l+m+mi}{0}\PY{p}{]}\PY{p}{,} \PY{n}{B}\PY{p}{[}\PY{l+m+mi}{1}\PY{p}{]}\PY{p}{)} \PY{c+c1}{\PYZsh{} Bandbreite mit Unsicherheit statistisch}

\PY{n}{k} \PY{o}{=} \PY{n}{m} \PY{o}{/} \PY{p}{(}\PY{l+m+mi}{4} \PY{o}{*} \PY{n}{T} \PY{o}{*} \PY{n}{B\PYZus{}ufloat}\PY{p}{)}

\PY{n+nb}{print}\PY{p}{(}\PY{l+s+sa}{f}\PY{l+s+s1}{\PYZsq{}}\PY{l+s+s1}{Boltzmannkonstante: k\PYZus{}B = }\PY{l+s+si}{\PYZob{}}\PY{n}{k}\PY{l+s+si}{:}\PY{l+s+s1}{.3uP}\PY{l+s+si}{\PYZcb{}}\PY{l+s+s1}{ J/K}\PY{l+s+s1}{\PYZsq{}}\PY{p}{)}

\PY{n+nb}{print}\PY{p}{(}\PY{l+s+sa}{f}\PY{l+s+s1}{\PYZsq{}}\PY{l+s+s1}{systematischer Fehler: }\PY{l+s+si}{\PYZob{}}\PY{n}{k}\PY{o}{.}\PY{n}{nominal\PYZus{}value}\PY{+w}{ }\PY{o}{*}\PY{+w}{ }\PY{l+m+mf}{0.02}\PY{l+s+si}{:}\PY{l+s+s1}{.3e}\PY{l+s+si}{\PYZcb{}}\PY{l+s+s1}{ J/K}\PY{l+s+s1}{\PYZsq{}}\PY{p}{)}
\end{Verbatim}
\end{tcolorbox}

    \begin{Verbatim}[commandchars=\\\{\}]
Boltzmannkonstante: k\_B = (1.39765±0.00392)×10⁻²³ J/K
systematischer Fehler: 2.795e-25 J/K
    \end{Verbatim}

    \begin{tcolorbox}[breakable, size=fbox, boxrule=1pt, pad at break*=1mm,colback=cellbackground, colframe=cellborder]
\prompt{In}{incolor}{57}{\boxspacing}
\begin{Verbatim}[commandchars=\\\{\}]
\PY{k+kn}{from}\PY{+w}{ }\PY{n+nn}{scipy}\PY{n+nn}{.}\PY{n+nn}{constants}\PY{+w}{ }\PY{k+kn}{import} \PY{n}{Boltzmann}

\PY{n}{sigma\PYZus{}dev} \PY{o}{=} \PY{p}{(}\PY{n}{k}\PY{o}{.}\PY{n}{nominal\PYZus{}value} \PY{o}{\PYZhy{}} \PY{n}{Boltzmann}\PY{p}{)}\PY{o}{/}\PY{n}{np}\PY{o}{.}\PY{n}{sqrt}\PY{p}{(}\PY{n}{k}\PY{o}{.}\PY{n}{std\PYZus{}dev}\PY{o}{*}\PY{o}{*}\PY{l+m+mi}{2} \PY{o}{+} \PY{p}{(}\PY{n}{k}\PY{o}{.}\PY{n}{nominal\PYZus{}value} \PY{o}{*} \PY{l+m+mf}{0.02}\PY{p}{)}\PY{o}{*}\PY{o}{*}\PY{l+m+mi}{2}\PY{p}{)}
\PY{n}{sigma\PYZus{}dev}
\end{Verbatim}
\end{tcolorbox}

            \begin{tcolorbox}[breakable, size=fbox, boxrule=.5pt, pad at break*=1mm, opacityfill=0]
\prompt{Out}{outcolor}{57}{\boxspacing}
\begin{Verbatim}[commandchars=\\\{\}]
np.float64(0.6022138763082926)
\end{Verbatim}
\end{tcolorbox}
        
    \hypertarget{ergebnis}{%
\subsection{Ergebnis}\label{ergebnis}}

Die berechnete Boltzmannkonstante beträgt
\[ k_B = (1.398\,\pm0.004\,\text{stat.}\,\pm 0.028\,\text{syst.}) \cdot 10^{-23} \frac{\text{J}}{\text{K}}\]

Die Sigma-Abweichung berechnet sich durch
\[ \sigma = \frac{k_B - k_{B,\text{lit}}}{\sqrt{(\Delta \text{stat})^2 + (\Delta \text{syst.})^2}} = 0.6 \Rightarrow 0.6\sigma \]

Dies spricht für eine sehr gute Übereinstimmung

Die getrennte Angabe von statistischem und systematischem Fehler ist
sinnvoll, da beide unterschiedliche Ursachen haben: Der statistische
Fehler resultiert aus den zufälligen Schwankungen des thermischen
Rauschens sowie der Einzelmessungen an den verschiedenen Widerständen
und lässt sich durch eine größere Anzahl an Messungen reduzieren. Der
systematische Fehler dominiert durch die 2\% Unsicherheit der
Rauschbandbreite. \(\Delta B\) bleibt für egal wieviele Messungen
gleich. Damit verschiebt wird Gesamtergebnis in eine feste Richtung
verschoben und kann nicht durch häufigeres Messen verringert werden.

Es zeigt sich durch die Trennung ausserdem, dass bereits genug Messungen
durchgeführt wurden, da der statitische Fehler um Faktor 10 kleiner als
der systematische Fehler ist.


    % Add a bibliography block to the postdoc
    
    
    

    
